\chapter{Theoretical background: Cartography} \label{theoretical background}

This chapter gives a general introduction to the theoretical background and serves as a reference to what my assumptions are based on and which amendments need to be made in Chapter \Ref{theoreticalassumptions}. The framework this book is written in is \textit{Cartography}, although it might be better to speak of a \textit{research field} or a ``\textit{research topic}" \citep[65]{cinriz:2010:carto}. Most assumptions of this book come from works by Guglielmo Cinque \citep{cinque:1993:evidence, cinque:2010:adj, cinque:2017:micro, cinque:2020:relcl, cinque:2021:carto, cinque:2023:linearization}, but I also incorporate ideas from \citet{laenzlinger:2005:french, laenzlinger:2011:elements, laenzlinger:2015:the, svenonius:2008:position} and \citet{kim:2019:syntax}. I start by outlining the core ideas of Cartography, before turning to direct and indirect modification as given by \citet{cinque:2010:adj}. I then discuss the additional approaches and close with an overview of previous studies of Japanese in the cartographic program.

\section{Basic ideas and concepts} \label{sectioncarto}
\largerpage[-3]
What Cartography offers is an elaborately articulated structure, a \textit{map} staying true to its name, of the noun phrase and the clause valid in every language. This research program is an offspring of X'-modeling and its adaptation to functional heads \citep{chomsky:1986:barriers} and developed parallel to Minimalism. Following Pollock’s fine-grained universal structure of the TP area (published in \citealt{pollock:1989:verb}), the clausal spine was constantly refined, most noteworthy through Rizzi’s split-CP-hypothesis \citep{rizzi:1997:the} and the highly articulated hierarchy of adverbials across different levels of the clause in \citet{cinque:1999:adverbs}. Following the doctoral thesis of Steven \citet{abney:1987:the}, in the late 1990s and early 2000s the noun phrase was extended with several projections as well, including projections for demonstratives \citep{bruge:2002:positions}, numerals \citep{ritter:1991:two, ritter:1993:where} as well as left-peripheral focus and topic projections \citep{dimitrova:1998:fragments, ihsane:2001:specific, aboh:2004:topic, aboh:2006:when} and of course projections for (adjectival) modifiers, the focal point of this book.\footnote{See \citet{alex:2007:nounphrase} for a large overview of the noun phrase, evidence for these projections and a historical overview of the literature.} One goal of this endeavor, and Cartography in general, is to argue for a parallel structure of the clause and the noun phrase.\footnote{See \citet{cinque:1993:evidence, cinque:1999:adverbs, cinque:2010:adj, cinque:2020:relcl, rizzi:1997:the, rizzi:2004:the, cinriz:2010:carto} and for a list of publications until 2015 \citet[144–151]{cinriz:2016:functional}. As also becomes obvious there, the clause has been the locus of much more attention, and far less controversy.}

The essential idea of Cartography is that the extended projection of each lexical head, N and V, contains several functional heads (F) where features are encoded and which give merging and landing sites in their specifier position for lexical and functional elements. Crucially, every functional projection is specified for which element(s) it hosts. Furthermore, these functional heads are ordered, are mapped, in a hierarchy, a map, that is based on structural selection and is argued to be universal across languages. This hypothesis also entails that every functional category attested in a given language is present in the structural make-up of all languages. As of now, we have no way of knowing how many functional projections there actually are in the clause and the noun phrase \citep[154–157]{cinriz:2016:functional}.

Languages thus essentially show parametric differences only in the lexicon, not in the syntax. They differ in how (and if) they realize each functional head and its specifier in the surface structure, as well as the type of movement they admit, if any. Essentially, this has the consequence that the structure offered by Cartography, although rigid at first glance and in a way the ``strongest position linguists can take" \citep[9]{bortolotto:2016:RA}, is constantly the subject of development and refinement, with every new language and micro-parameter investigated, as well as every new approach and perspective taken. This makes this research field not only highly relevant and powerful, but also extremely versatile and subject to constant change.\footnote{See \citet{cinque:1993:evidence, cinque:1999:adverbs, cinque:2010:adj, cinque:2020:relcl, rizzi:1997:the, rizzi:2004:the, shlonsky:2004:form, poletto:2000:the, scott:2002:stacked, laenzlinger:2011:elements} and many more. For an overview, see \citet{cinriz:2010:carto, cinriz:2016:functional}.}

\subsection{Ordering}

The rigid order of functional projections (correctly) predicts the existence of ordering restrictions of the type we saw in English in the introduction. A well-established \is{adjective ordering restrictions}ordering hierarchy of adjectives attested in several unrelated languages \citep{sussex:1974:the, dixon:1982:where, hetzron:1978:on} and partially attributed to cognitive factors \citep{sproat:1987:pre, sproat:1991:the} reads as follows, where the symbol `$\succ$’ means `\textit{precedes in linear order}’.

\ex. \textsc{quality} $\succ$ \textsc{size} $\succ$ \textsc{shape} $\succ$ \textsc{color} $\succ$ \textsc{provenance} $\succ$ N \\ \null \hfill \citep{sproat:1991:the}

It was a central achievement of the cartographic research field (most notably \citealt{cinque:1993:evidence, scott:2002:stacked, laenzlinger:2005:french, panayidou:2013:in} and \citealt{cinque:2010:adj}) to attribute these ordering principles to structural selection, getting rid of adjunction, which cannot explain ordering preferences.

In other words, each of these semantic categories corresponds to a functional head and adjectives are merged as specifiers of the head that matches their semantic value. An adjective denoting \textsc{size} would merge in FP\textsubscript{\textsc{size}} and so on.

The hierarchical \is{adjective ordering restrictions}order of these heads can be represented as functional selection in a bracketing format \ref{bracketing} or as in \ref{hierarchy no bracketing} where the symbol `$>$' means `\textit{is structurally/hierarchically higher}’.\footnote{When an element A is hierarchically higher than an element B, it precedes it in structural, but not necessarily in linear order since movement operations might yet take place. The symbol `$\succ$’ means `precedes in linear order’. The same technique is applied in \citet{panayidou:2013:in}.}

\ex. [\textsc{quality} [\textsc{size} [\textsc{shape} [\textsc{color} [\textsc{provenance} [N ]]]]]] \label{bracketing}

\ex. \textsc{quality} $>$ \textsc{size} $>$ \textsc{shape} $>$ \textsc{color} $>$ \textsc{provenance} $>$ N \label{hierarchy no bracketing}

This functional hierarchy is expected to be cross-linguistically valid and was verified among others for Romance, \is{Greek} Greek, \is{Romanian} Romanian, \is{Welsh} Welsh, \is{Mandarin} Mandarin (see for example \citealt{scott:2002:stacked}, the appendix of \citealt{cinque:2010:adj} and many examples in \citealt{alex:2007:nounphrase}).

\subsection{Left-right asymmetries} \label{left-right}

A fundamental contribution of Cartography, first systematized in \citet{cinque:2005:deriving, cinque:2009:fundamental} and extended in \citet{cinque:2023:linearization, cinque:2025:superiority}, was to shed light on so-called \textit{Left-right asymmetries}. It has been a central cross-linguistic observation (see for example \citealt{rijkhoff:2002:noun}), and Greenberg's Universal 20 \citep{greenberg:1963:some}, that there is a restriction on possible orders of modifiers around a lexical head when they are all on the same side of the head. When they appear on the left side of the lexical head, they are ordered in only one specific way, but when they follow the lexical head, two orders are possible: Either the identical order maintaining linear precedence, or the mirror-image, in which the order has been reversed compared to the pre-head order. This observation was verified for the order of adjectives \citep{cinque:2009:fundamental, shlonsky:2004:form}, but for many other domains as well, including the order of adverbs as well as circumstantial and directional/locative prepositions around the verb \citep[167–168]{cinque:2009:fundamental}. \is{left-right asymmetries} \is{U20}

A very illustrative case, and perhaps the central case of application in the nominal domain, is the order of demonstratives, numerals and adjectives with respect to the noun. Whereas languages such as English or German reflect the underlying order \textit{demonstrative} $\succ$ \textit{numeral} $\succ$ \textit{adjective} $\succ$ \textit{noun}, in languages in which these constituents appear postnominally (in canonical order), two possibilities are found: The relative order of modifiers as they appear before the noun in English, or the mirror-image of the English order. Crucially, however, the mirror-image order in prenominal position is not attested as the natural order of constituents in any language. This is visualized below.

\ex. Order of demonstratives, numerals and adjectives \citep[166]{cinque:2009:fundamental}
\a. Dem $\succ$ Num $\succ$ A $\succ$ N (English, German...)
\b. N $\succ$ Dem $\succ$ Num $\succ$ A (Abu', Kikuyu...)
\c. N $\succ$ A $\succ$ Num $\succ$ Dem (Gungbe, Thai...)
\d. *A $\succ$ Num $\succ$ Dem $\succ$ N (not attested)

Although humans do have a strong bias towards homomorphism, that is, placing all elements on the same side of a lexical head (cf. \citealt{martin:2020:experimental}, references therein and further work by the same authors), in many different languages, modifiers appear on both sides of a lexical head in linear order. The crucial point is however that not all mathematically possible 24 (4 factorial) orders are attested as base orders across the world’s languages, but only 14 out of 24 (\citealt{cinque:2005:deriving, cinque:2023:linearization}, and see already \citealt{greenberg:1963:some, hawkins:1983:word}). Besides the order A $\succ$ Num $\succ$ Dem $\succ$ N, neither Num $\succ$ Dem $\succ$ N $\succ$ A nor Dem $\succ$ A $\succ$ Num $\succ$ N is attested \citep[319–320]{cinque:2005:deriving}. \is{left-right asymmetries} \is{U20}

Under the anti-symmetric approach to linearization of Kayne’s \citeyear{kayne:1994:anti} \is{LCA}Linear Correspondence Axiom (LCA), according to which there is only one structural hierarchy of elements,\footnote{``d(\textit{A}) is a linear ordering of \textit{T}" \citep[6]{kayne:1994:anti}.} and the relation between these elements is governed via \textit{asymmetric c-command},\footnote{``Let X, Y be nonterminals and x, y terminals such that X dominates x and Y dominates y. Then if X asymmetrically c-commands Y, x precedes Y" \citep[31]{kayne:1994:anti}; see also \citet{moro:1997:raising, moro:2000:dynamic}.} those left-right asymmetries have been argued to fall out naturally, given the assumption that only one underlying hierarchy exists and all other orders arise via movement. The purpose of movement in Cartography is thus primarily the alteration of linear precedence, not feature checking as in minimalist frameworks.

Second, one needs to take into account that movement is restricted, otherwise we would see all of the 24 possible orders attested cross-linguistically. It is only the lexical head that moves and only via phrasal movement to the left, but in different fashions. It moves either alone, crossing each higher functional projection it is c-commanded by (and possibly halting in-between causing non-homomorphic orders), or in a \is{pied-piping}pied-piping fashion (also called \textit{roll-up} or \textit{snowballing}, \citealt{ross:1967:constraints}) where it takes each constituent present in each functional projection it crosses along and reverses their order \citep{cinque:2017:micro, cinque:2023:linearization}. When the noun has reached its destination, the order of all elements it has taken along has been reversed accounting, for instance, for the order N $\succ$ A $\succ$ Num $\succ$ Dem \citep[318]{cinque:2005:deriving}. Pied-piping can also occur partially, deriving for example the order Dem $\succ$ Num $\succ$ N $\succ$ A in which the NP has moved to an agreement position above the adjective. Orders such as *Num $\succ$ Dem $\succ$ N $\succ$ A can, however, never be achieved since the order Num $\succ$ Dem can only be achieved if the noun moves across these categories, which it has not done \citep[317, 318]{cinque:2005:deriving}.\footnote{Matters are more complex. There is another type of pied-piping, namely the so-called \textit{pictures-of-whom} or \textit{regressive} pied-piping, in which the head moves and takes its constituents with it, however \textit{without} reversing the order. Furthermore, several movement operations of different types can combine in one derivation. An important argument for the complexity of these movement operations is the derivation of certain rare constituent orders across languages. For instance, if we take Dem $\succ$ Num $\succ$ A $\succ$ N as the base order, the order Num $\succ$ A $\succ$ N $\succ$ Dem, for which \citet{cinque:2005:deriving, cinque:2023:linearization} finds 51 languages, can only be derived if the NP moves above the functional projection hosting NumP in a \textit{pictures-of-whom} pied-piping fashion, not reversing the order, and then continues moving above DemP in a \textit{whose-pictures} type fashion, this time reversing the order \citep[36]{cinque:2023:linearization}. See \citet[9–11]{cinque:2025:superiority} for an application to adjective order. For a recent contribution on U20-effects in a nanosyntactic framework see \citet{pinzin:2025:word}. \label{footnotemovement}} \is{U20}

These mechanics can be extended to all other domains mentioned above as well, that is the order of adjectives, adverbs and prepositions, among others. As recently argued in \citet{cinque:2025:superiority}, it is exactly this point, the applicability across different domains and the prediction on which orders are attestable, that constitutes a big advantage of the movement approach over other approaches. I will include a brief discussion on how the account proposed in this book fares with these observations in Section \ref{discussion u20}.

\subsection{The Japanese DP: No NP-movement} \label{japanese no movement}

Admittedly, there is a skew in cartographic research towards languages of the Indo-European family. Japanese is considered in \citet[3–4]{cinque:2017:micro} a ``closest to the ideal'' head-final language. This is because modifiers (demonstratives, relative clauses, numerals, adjectives) systematically precede the noun, but most functional morphemes (suffixes for plurality or case) follow it.\footnote{Exceptions are honorific morphemes, which are realized as prefixes.} An example of a Japanese DP containing an adjective and a plural suffix is illustrated in \ref{wakai}.

\exg. [Waka-i gakusei-tachi-ga] ki-ta.\\
young-\textsc{i} student-\textsc{pl-nom} come-\textsc{pst}\\
`The young students came.’ \label{wakai}

As modifiers always precede the noun, no NP-movement needs to be assumed in the strict sense. Under the assumption that adjectives such as \textit{waka-i} `young’ are embedded in specifier positions of functional projections, but number suffixes are heads, order falls out naturally as depicted in the structural representation of \ref{wakai} in \figref{tree wakai}. Note that I deviate here from the strict application on Kayne’s \is{LCA}LCA in depicting Japanese trees as head-final. Movement does become necessary if we were to depict Japanese structures as head-initial.\footnote{This point also extends to other languages and phenomena. See \citet{buering:1997:doing} for a discussion on German extraposed relative clauses, which are hard to analyze with the rigid assumptions of the LCA. See \citet{abels:2012:.linear} for general discussion of the problems this axiom entails.} \is{German}

\begin{figure}
\centering
\begin{tikzpicture}[baseline=0pt, remember picture]
\tikzset{level distance=25pt, sibling distance=5pt}
 \Tree [.DP {} [.D′ [.AgrP {} [.Agr′ [.NumberP {} [.Number′ [.AgrP {} [.Agr′ [.FP [.AP \edge[roof]; waka-i ] [.F′ [.NP \edge[roof]; gakusei ] F ] ] Agr ] ] [.Number [.Pl tachi ] ] ] ] Agr ] ] D ] ]
\end{tikzpicture}
\caption{Structural representation of \ref{wakai}}
\label{tree wakai}
\end{figure}

\section{Direct and indirect modification} \label{cartography Cinque}

In this section, I review the differences between direct and indirect modifiers, first, in order to derive reliable diagnostics for the type of modification, second, to relate their different structural positions to these diagnostics. In the following two sections, I elaborate on direct-only modifiers and the derivation of indirect modifiers. Discussions will be kept short here as they will be picked up in the following chapters.

\subsection{Introduction}

The notions \textit{direct} and \textit{indirect modification} were established by \citet{sproat:1991:the}. The authors used them as an explanation for why adjectives in \is{Mandarin} Mandarin Chinese are ordered rigidly in some instances, but not in others. More precisely, they noted that an overt exponent of indirect modification is the element \textsc{de} \textjpn{的} \citep[§ 2]{sproat:1991:the}.\footnote{Analyzing the precise nature of \textsc{de} has been the task for a long time in syntactic research. \citet{sproat:1991:the} analyze it as a complementizer (see also \citealt{gobbo:2010:on}), but authors including Cinque (\citeyear{cinque:2010:adj}) do not share this proposal and have shown that some direct modifiers in fact have to appear with this element, (cf. \citealt{paul:2005:adj, paul:2010:adj, paul:2012:why, yang:2005:plurality}) and analyze it as `phrasal modifier'. Other analyses include `case assigner' \citep{larson:2009:chinese} or `linker' \citep{dikken:2004:complex}. \label{demandarin}} This element is optional in adjectival modification, but when it appears, it causes the absence of otherwise attested \is{adjective ordering restrictions}ordering restrictions, which are otherwise similar to the English case \citep[579]{sproat:1991:the}. This is shown below.

\ex. \textsc{size} $\succ$ \textsc{color} without \textsc{de}
\ag. xi\v{a}o lù huāpíng\\
small green vase\\
\bg. ?lù xia\v{o} huāpíng\\
green small vase\\
`small green vase’ \null \hfill \citep[566]{sproat:1991:the}

\ex. \label{mandarin with de} \textsc{size} vs. \textsc{color} with \textsc{de}
\ag. xi\v{a}o de lù de huāpíng\\
small \textsc{de} green \textsc{de} vase\\
\bg. lù de xi\v{a}o de huāpíng\\
green \textsc{de} small \textsc{de} vase\\

Furthermore, this element is mandatory in verbal modification, which arguably always involves a relative clause structure.

\exg. [nèizhī jiào *(de)] g\v{o}u\\
that bark \textsc{de} dog\\
`the dog which is barking' \null \hfill \citep[573]{sproat:1991:the}

As mentioned above, Cartography predicts that direct adjectives are merged in the specifier positions of functional heads that match their semantics and which are rigidly ordered. And since FP\textsubscript{\textsc{size}} structurally dominates FP\textsubscript{\textsc{color}}, linear adjacency is expected in context-neutral utterances. The hierarchical structure is depicted below for English \citep[28–30]{cinque:2010:adj}.\footnote{The threshold between the direct and the indirect domain is the indefinite head d. See Section \ref{indefinitedpdiscussion}.} \is{English} \is{Mandarin} \is{adjective ordering restrictions}

\ex. \a. small green vase
\b. ??green small vase
\c. \begin{tikzpicture}[baseline=0pt, remember picture]
\tikzset{level distance=25pt, sibling distance=5pt}
\Tree [.DP {} [.FP\textsubscript{indirect} XP [.dp d [.FP\textsubscript{\textsc{size}} [.AP \edge[roof]; {small} ] [.F′\textsubscript{\textsc{size}} F [.FP\textsubscript{\textsc{color}} [.AP \edge[roof]; {green} ] [.F′\textsubscript{\textsc{color}} F [.NP \edge[roof]; \node (voitures4) {vase}; ]]]]]]]]]
\end{tikzpicture}

Indirect modifiers, on the other hand, are merged in the projection of relative clauses and since these functional projections are unordered with respect to semantic category, it is expected that indirect modifiers exhibit free order on the surface. This is what is at stake in Mandarin in \ref{mandarin with de}, claimed \citet{sproat:1991:the}, and this behavior is also observable in English (\citealt{baker:2003:verbal}).\footnote{Speakers might disagree on the possibility of relative clause stacking. And while \ref{bigredrc} does sound a bit stilted and was chosen mostly for the parallelism to the Mandarin example, the same as the German example in \ref{bigredrcgerman}, examples of stacked relative clauses can be found, although preferably with longer relative clauses such as ``\textit{a book [that I wanted to read for a long time] [that was recommended to me by my teacher]}".}

\ex. \label{bigredrc} \a. a vase that is small that is green
\b. a vase that is green that is small

\subsection{Reduced relative clauses and ambiguities} \label{ambiguities}

As mentioned in the introduction, most adjectives in English are actually ambiguous between a direct and an indirect status. Direct modifiers, here called \textit{direct-only}, are a very rare species. In fact, all adjectives that can have access to predicative use can have an indirect status when they appear in a simple modification structure. They are then not embedded under a full (also: finite) relative clause, but are instances of so-called \is{reduced relative clause}\textit{reduced relative clauses} (RRC) (\citealt[54–55]{cinque:2010:adj}, \citealt{douglas:2016:syntactic, harwood:2018:rrc}), relative clauses which lack a relative pronoun/complementizer and a verb/copula. This fact has also been put forth to explain deviations from expected ordering restrictions. If two indirect adjectives modify a noun, no fixed order is predicted, and languages where adjectives never follow ordering restrictions are predicted to only have indirect modification adjectives (cf. \citealt[§ 3.3.2]{cinque:2010:adj}).

Verbal reduced relative clauses are easily recognizable in \is{English} English, namely by virtue of lacking relative pronouns and copula, as for instance is the case with participial and infinitival relative clauses \citep{kayne:1994:anti, cinque:2020:relcl}, but there are also clear cases of indirect adjectives embedded in reduced relative clauses. Prototypical examples are postnominal adjectives of the type \textit{the man proud of his children} \citep{williams:1982:fofc, kayne:1994:anti, sadler:1994:prenominal, larson:2000:acd, larson:2004:indefinite}. While most adjectives in English are only licit in this position when they occur with complements or adjuncts,\footnote{The \textit{Head-Final Filter} of \citet{williams:1982:fofc}, cf. \citet{alexeyenko:2021:hff, richards:2022:finding, richards:2023:finding}.} there are some adjectives that can appear independently in this position, in addition to the prenominal position. Those include adjectives with the ending -\textit{ible/able} as in the following famous example from \citet[3--4]{bolinger:1967:adj}.\footnote{Adjectives with the initial prefix \textit{a}- such as \textit{afraid} or \textit{asleep} predominantly occur in postnominal position (\citealt{bolinger:1967:adj, larson:1998:events, larson:2004:indefinite}; \citealt[295]{alex:2007:nounphrase}; \citealt[§ 2.3]{belk:2017:attributes}).}\is{Head-Final Filter}

\ex. \a. the (only) visible stars
\b. the (only) stars visible

It is well-known that such minimal pairs in English exhibit a meaning difference. In postnominal position, the adjectives denote a temporary, passing property, dubbed \is{stage/individual-level}\textit{stage-level} reading. In prenominal position, they are ambiguous and can have this reading but additionally can also be understood as denoting an enduring property dubbed \textit{individual-level} reading (see \citealt{sadler:1994:prenominal, larson:1998:events, larson:2004:indefinite} and especially \citealt[§ 2.1]{cinque:2010:adj}). This allows the following paraphrases.

\ex. \a. the (only) stars visible $\approx$ `the only stars visible at the moment/that happen to be visible' (stage-level)
\b. the (only) visible stars $\approx$ `the only stars visible at the moment/that happen to be visible' (stage-level) OR `the only stars visible in general' (individual-level)

There are many such minimal pairs in English. What is common to each of them is that it is always the prenominal variant that is ambiguous, while the postnominal variant is restricted to a certain reading as shown below for the famous case of \textit{beautiful dancer} which can either have an intersective or a non-intersective reading (\citealt{larson:2000:acd, vendler:1957:verbs, siegel:1976:capturing, larson:1995:olga}). \is{intersectivity}

\ex. \label{ambiguous, beautiful} Non-Intersective vs. Intersective \citep[145]{larson:1995:olga}
\a. Olga is a beautiful dancer. (ambiguous) \label{olga}
\b. $\approx$ `Olga is a dancer that dances beautifully.’ (non-intersective reading)
\c. $\approx$ `Olga is a dancer and Olga is beautiful.’ (intersective reading)

\ex. Non-Intersective vs. Intersective \citep[9]{cinque:2010:adj}
\a. Olga is a dancer more beautiful than her instructor. (only intersective)
\b. $\neq$ `Olga is a dancer that dances beautifully.’ (non-intersective reading)
\c. $\approx$ `Olga is a dancer and Olga is beautiful.’ (intersective reading) \label{intersective, RC}

Cartography offers an explanation for this: The ambiguity in interpretation observed in prenominal position arises since the adjective is syntactically ambiguous - it could be either direct or indirect. On the other hand, in postnominal position, the adjective can only be indirect. This can be verified when embedding it in a full relative clause. Again, only one reading is available: The same as in the postnominal use. This confirms that postnominal adjectives in \is{English} English are predicative in nature \citep[18–19]{cinque:2010:adj}.

\ex. Non-Intersective vs. Intersective, full RC \citep[9]{cinque:2010:adj}
\a. Olga is a dancer who is beautiful. (only intersective)
\b. $\neq$ `Olga is a dancer that dances beautifully.’ (non-intersective reading)
\c. $\approx$ `Olga is a dancer and Olga is beautiful.’ (intersective reading)

As a consequence, the reading not accessible in postnominal position, and not derived from predicative use, must be the direct reading and when adjectives such as \textit{beautiful} exhibit that reading, they must be interpreted as direct.

\subsection{Direct-only modifiers}

There is a small group of adjectives that are unambiguously direct. Those are called \textit{direct-only} in this book. According to \citet[§ 4.1]{cinque:2010:adj}, these modifiers are \textit{functional} and \textit{phrasal}. They are functional because they arguably form a closed class, in most languages only with a couple of members \citep[43–44]{cinque:2010:adj}, explaining many cross-linguistic intersections. They are phrasal because they can appear with complements and have a phrasal character in some languages \citep[44–49]{cinque:2010:adj}. Below, I will review some well-known members as well as characteristics, but it should be kept in mind that both are, most of all, prototypical.

A well-known group of direct-only modifiers is \is{non-intersective modifier}\textit{non-intersective}, also called \textit{intensional}, adjectives \citep{kamp:1995:prototype, panayidou:2013:in}, such as \textit{former}. Their inability to appear in predicative position, and thus also inside a relative clause, has been used to refute earlier transformational accounts à la \citet{smith:1964:det} and is a key diagnostic for direct-only status.

\ex. \a. a former president
\b. *This president is former.
\c. *a president who is former

In contrast to the \textit{beautiful dancer} example above, such adjectives only yield a non-intersective interpretation, also in attributive position, and are not ambiguous.

\ex. \a. Marya is a former dancer. (only non-intersective)
\b. $\approx$ `Marya was formerly a dancer.’ (non-intersective reading)
\c. $\neq$ `Marya is former, and Marya is a dancer.’ (intersective reading) \\ \null\hfill \citep[304]{alex:2007:nounphrase}

Further typical characteristics of direct-only modifiers include the inability to co-occur with degree expressions (Section \ref{gradability}) and to exhibit overt tense (Section \ref{test tense}).

\ex. *the very former president

\ex. *the (once) (was) former president / president who was former

Another \is{non-intersective modifier}type of direct-only modifier occurs in \is{idiomatic modifier}idiomatic expressions. Such modifiers, those participating in idiomatic constructions, naturally diverge across languages. In English, they must occur prenominally. They can neither occur postnominally, nor in predicative position and still keep their idiomatic meaning \citep[§ 7.2]{cinque:2010:adj}.

\ex. \a. a white lie \hfill \citep[88]{cinque:2010:adj}
\b. *a lie white (in spirit)
\c. *a lie (that) is white

\subsection{Linear distance}

Indirect modifiers are structurally higher than direct modifiers, which predicts linear adjacency of direct modifiers to the noun. When an ambiguous modifier occurs twice in prenominal position in \is{English} English, it does so once with each reading \citep[19–20]{cinque:2010:adj}. Although the intuition is subtle, the modifier closest to the noun clearly seems to exhibit the direct reading as shown in \ref{visiblevisible} (cf. \citealt[280–281]{larson:2004:indefinite}).

\exg. The visible visible stars include Capella. \hfill \citep[155]{larson:1998:events}\\
the stage-level individual-level stars\\ \label{visiblevisible}

The two modifiers can also be spread out on both sides of the noun. Then, as expected, the postnominal modifier must bear the indirect modification, the \is{stage/individual-level}stage-level reading \citep[19–20]{cinque:2010:adj}.

\exg. The visible stars visible include Capella. \null\hfill \citep[155]{larson:1998:events}\\
the individual-level stars stage-level\\

These observations hold for the other dichotomies, for example the intersective/non-intersective case. \is{intersectivity}

\largerpage[1]

\ex. \ag. Olga is a beautiful beautiful dancer.\\
Olga is a intersective non-intersective dancer\\
\null \hfill \citep[281]{larson:2004:indefinite}
\bg. Olga is a beautiful dancer more {beautiful than her instructor.}\hspace{-5mm}\\
Olga is a non-intersective dancer more intersective\\
~\hfill\citep[19]{cinque:2010:adj}

\subsection{More on ordering restrictions} \label{AORs}

Standard Cartography explains \is{adjective ordering restrictions}ordering restrictions among direct modifiers as a result of the distinct positions of the respective functional heads in the direct domain. Several hierarchies of functional heads for direct modification have been proposed, but a standard hierarchy for object nouns could look as follows. I have chosen the label \textsc{dimension} as a more broad term than \textsc{size}. It also includes the categories \textsc{length} and \textsc{width} from \citet[114]{scott:2002:stacked}.\footnote{Modifiers of \textit{action} or \textit{event} nouns \citep{grimshaw:1991:argument}, correspond to adverbs in the verbal domain \citep{cinque:1993:evidence, cinque:1999:adverbs, alexiadou:1997:adverb} and are often given in the following hierarchy: \textsc{speaker-oriented} $>$ \textsc{subject-oriented} $>$ \textsc{manner} $>$ \textsc{thematic} $>$ N \citep[324]{alex:2007:nounphrase}.}

\ex. \hspace*{-2mm}\textsc{quality} $>$ \textsc{dimension} $>$ \textsc{shape} $>$ \textsc{color} $>$ \textsc{nationality} $>$ \textsc{material} $>$ N \label{hierarchyjo}

Importantly, ordering restrictions are only maintained in \textit{hierarchical} modification. In \textit{parallel} modification, achieved via syndetic or asyndetic \is{coordination}coordination, no ordering restrictions are expected \citep[578–580]{sproat:1991:the}. \is{English}

\ex. Parallel Modification (English)
\a. a red and big book \null \hfill syndetic
\b. a red, big book \null \hfill asyndetic

\subsection{Summary: Direct vs. indirect modifiers}
\largerpage[2.5]
The defining characteristics that differentiate direct from indirect modifiers which offer valid diagnostics and tests to determine the status of a given modifier are summarized in \tabref{table differences direct and indirect} (cf. \citealt[33]{cinque:2010:adj} for a more detailed version).

\begin{table}
\begin{tabularx}{\textwidth}{QQ}
\lsptoprule
\textbf{Direct} & \textbf{Indirect} \\
\midrule
 Exhibit special readings, such as non-intersective, modal etc. & Exhibit the opposite readings: intersective, implicit-RC etc. \\
\\[-.7em]
Take part in idiomatic constructions & Only give literal interpretation \\
\\[-.7em]
Impossible as predicates & Possible as predicates \\
\\[-.7em]
Closer to the head noun & Further away from the head noun \\
\\[-.7em]
Ordered rigidly & Show free/arbitrary ordering\\
\\[-.7em]
Tend to be more simple & Are usually complex\\
\lspbottomrule
\end{tabularx}
\caption{Differences between direct and indirect modifiers}
\label{table differences direct and indirect}
\end{table}

%\begin{table}
%\begin{tabular}{ll}
%\lsptoprule
%Direct & Indirect \\
%$\midrule
%Exhibit special readings, such as non-intersective, modal etc. & Exhibit the opposite %readings: intersective, implicit-RC etc. \\
%Take part in idiomatic constructions & Only give literal interpretation \\
%Impossible as predicates & Possible as predicates \\
%Closer to the head noun & Further away from the head noun \\
%Ordered rigidly & Show free/arbitrary ordering \\
%Tend to be more simple & Are usually complex \\
%\lspbottomrule
%\end{tabular}
%\caption{Differences between direct and indirect modifiers}
%\label{table differences direct and indirect}
%\end{table}

\subsection{Cross-linguistic perspective}

It is an important enterprise of the cartographic research programs to verify those facts in cross-linguistic perspective. And indeed, they can, and often do, serve as valid diagnostics for modification status in other languages as well. A cartographic account of the DP has been proposed for \is{French} French \citep{laenzlinger:2005:french, laenzlinger:2011:elements}, \is{Greek} Greek (\citealt[appendix]{cinque:2010:adj}; \citealt{panayidou:2013:in}), \is{Russian} Russian (\citealt{siegel:1976:capturing}; \citealt[appendix]{cinque:2010:adj}), \is{Dutch} Dutch \citep{belk:2017:attributes}, \is{Korean} Korean \citep{kang:2005:adj, kang:2006:two, kim:2019:syntax}, \is{Javanese} Javanese \citep{vander:2009:direct} and \is{Kipsigis} Kipsigis \citep{kouneli:2019:syntax}.\footnote{See also the appendix of \citet{cinque:2010:adj} for \is{Romanian} Romanian, \is{Maltese} Maltese and \is{German} German and the typological sample in \citet{laenzlinger:2010:cp, laenzlinger:2011:elements}.} That \is{adjective ordering restrictions}ordering restrictions, or at least tendencies, have cross-linguistic validity is exemplified for German in \ref{bigredgerman}. Again, once the modifiers are embedded in corresponding relative clauses, their order becomes free \ref{bigredrcgerman}.

\ex. \label{bigredgerman} \textsc{size} vs. \textsc{color} (German)
\ag. ein großes rotes Buch\\
a big red book\\
\bg. ??ein rotes großes Buch\\
a red big book\\
`a big red book’

\ex. \label{bigredrcgerman} \textsc{size} vs. \textsc{color} indirect (German)
\ag. ein Buch, das groß ist, das rot ist\\
a book that big is that red is\\
\bg. ein Buch, das rot ist, das groß ist\\
a book that red is that big is\\
`a book which is big which is red’

As argued by \citet{cinque:2010:adj}, the resolution of ambiguities also holds in \is{Italian} Italian (Romance), where adjectives are freer in their distribution to the left and to the right of the noun. Ambiguous adjectives of the \textit{beautiful}-type, arguably, show this ambiguity in their prototypical position as well, but this is the postnominal position \ref{italianattacantebuono}. On the contrary, in the more marked, prenominal position, the relevant adjectives only have one interpretation. This interpretation is the opposite of the English postnominal cases, thus corresponds to the direct reading \ref{italianbuonoattacante}. From this, \citet[§ 2.4]{cinque:2010:adj} concludes that prenominal adjectives in Italian must be direct \is{intersectivity} (examples taken from \citealt[10]{cinque:2010:adj}).

\exg. Un \textbf{attaccante} \textbf{buono} non farebbe mai una cosa del genere.\\
a forward good not would.do never a thing of.the kind\\
1.. $\approx$ `A forward good at playing forward would never do such a thing.’ (non-intersective)\\
2. $\approx$ `A good-hearted forward would never do such a thing.’ (intersective) \label{italianattacantebuono}

\exg. Un \textbf{buon} \textbf{attaccante} non farebbe mai una cosa del genere.\\
a good forward not would.do never a thing of.the kind\\
1. $\approx$ `A forward good at playing forward would never do such a thing.’ (non-intersective)\\
2. $\neq$ `A good-hearted forward would never do such a thing.’ (intersective)\label{italianbuonoattacante}

The same ambiguous adjective with two different readings is only possible postnominally. Nevertheless, the relative closeness is preserved: The adjective with the direct modification interpretation must be closer to the noun than the adjective with the indirect modification reading \citep[20–22]{cinque:2010:adj}. \is{intersectivity}

\exg. Un attaccante \textbf{buono} buono \\
a forward good(non-intersective) good-hearted(intersective)\\
`a good-hearted good forward' \null\hfill \citep[21]{cinque:2010:adj}

\section{Groups of direct-only modifiers} \label{directmodifiers}

This section presents the two biggest cross-linguistic groups of direct-only modifiers. Besides non-intersective modifiers, which have already been discussed, I will discuss \textit{relational} modifiers. Both these groups are relevant for the discussion of direct modifiers in Japanese in Chapter \ref{direct japanese}.

\subsection{Non-intersective modifiers} \label{intersective}

\subsubsection{Truly non-intersective modifiers}

Non-intersective modifiers of the type \textit{former} or \textit{alleged} are characterized by the fact that they do not denote a set intersection with the noun they modify \citep{vendler:1957:verbs, siegel:1976:capturing, bolinger:1967:adj, kamp:1995:prototype, larson:1998:events}. This is illustrated with the English \is{non-intersective modifier}adjective \textit{former} below where there is no set intersection between the set of presidents and the set of former entities pertaining to Obama.

\ex. \a. Obama is the/a former president.
\b. [former president] $\neq$ [former] $\cap$ [president]

As predicative use enforces the intersective reading, it is only logical that these modifiers cannot be used predicatively.

\ex. *This president is former.

Such \textit{truly non-intersective} modifiers often refer to a time (past, future, present), a place (\textit{dortig}; \is{German} German: `over there’) or a manner (possible, main), highlighting their connection to adverbs. This group has also been recognized in Japanese \citep{yamakido:2005:nature, ayano:2010:revisiting} and behaves alike. Non-intersective modifiers in Japanese cannot be used predicatively and do not denote intersections with the noun.

\ex. \ag. Tarō-ga shōrai-no purosakkāsenshu-da.\\
Taro-\textsc{nom} future-\textsc{no} pro.soccer.player-\textsc{cop}\\
`Taro is a future professional soccer player.’
\bg. *Ano purosakkāsenshu-ga shōrai-da.\\
that pro.soccer.player-\textsc{nom} future-\textsc{cop}\\
(`That professional soccer player is future.’)
\c. [shōrai-no purosakkāsenshu] $\neq$ [shōrai] $\cap$ [purosakkāsenshu]

\subsubsection{Subsective modifiers}

Notably, there is an in-between class of modifiers between intersective and non-intersective modifiers, usually called \textit{subsective}, where modifiers co-occur with nouns depicting an event or a profession and only modify part of the extension of those nouns. A well-known example is the following \citep[137–138]{kamp:1995:prototype}.

\ex. Mary is a skillful surgeon.

The most salient interpretation of this phrase is that Mary is \is{intersectivity}skillful as a surgeon, so she is experienced in her profession of operating people. The interpretation `Mary is skillful’ again seems not really available, because we do not know what she could be skillful at. Subsective modifiers are able to appear in predicative position and thus also in relative clauses, but then only one interpretation remains, namely the intersective one. The subsective interpretation, on the other hand, is inaccessible.

\ex. \a. Mary is a surgeon who is skillful. (only intersective)
\b. $\neq$ `Mary is skillful as a surgeon.’ (subsective)
\c. $\approx$ `Mary is a surgeon and Mary is skillful.’ (intersective)

Subsectivity is the source of most ambiguities mentioned above, including the famous \textit{beautiful dancer} example (\citealt[145]{larson:1995:olga}; \citealt{vendler:1957:verbs, siegel:1976:capturing}). The paraphrase ``beautiful as a dancer" is in fact the subsective readig, while the paraphrase ``beautiful and a dancer" is the intersective reading.

\ex. \a. Olga is a beautiful dancer.
\b. [beautiful dancer] $\subseteq$ [dancer] (subsective)
\c. [beautiful dancer] = [beautiful] $\cap$ [dancer] (intersective)

Another well-known ambiguous example is the following.

\ex. \a. Peter is an old friend. \hfill \citep[146]{larson:1998:events}
\b. [old friend] $\subseteq$ [friend] (subsective: `long-time friend’)
\c. [old friend] = [old] $\cap$ [friend] (intersective: `aged friend’)

In Section \Ref{non-intersective Japanese}, it will be shown that Japanese nominal modification does not readily allow such ambiguities, one reason why direct modification has been denied. Instead, ambiguities are resolved via different lexemes, each exhibiting an intersective and a non-intersective reading.

\subsubsection{Privative modifiers}

Another group of non-intersective modifiers is \is{privative modifier}privative modifiers. Such modifiers, often similar in spirit to the meaning `fake’, denote a zero intersection with their head noun.

\ex. \a. a fake gun
\b. [fake gun] $\cap$ [gun] = $\emptyset$ \hfill \citep[25]{morzycki:2016:mod}

Interestingly, such modifiers can appear in predicative position, which is why \citet{partee:2009:formal} has proposed that \is{privative modifier}privative adjectives in \is{English}English and \is{Russian} Russian might be subsective after all, therefore not direct-only.

\ex. This gun is fake / pretend / artificial. \hfill \citep[25]{morzycki:2016:mod}

This is also the case in Japanese where predicative use of the privative modifier \textit{nise-no} `fake' is degraded, but not entirely ungrammatical.

\exg. ?Kono pisutoru-wa nise-da.\\
this gun-\textsc{top} fake-\textsc{cop}\\
`This gun is fake.'

As shown in \citet[21–23]{cinque:2014:semantics} however, there is reason to analyze privative modifiers as direct-only. Consider the following minimal pair with the \is{Italian} Italian equivalent \textit{falso} `false’.

\ex. \ag. le false banconote (con cui giocavano)\\
\textsc{det} fake bank.note.\textsc{pl} with which they.played\\
`the fake money with which they used to play’ (only: `Monopoly money’)
 \bg. le banconote false (con cui giocavano)\\
\textsc{det} bank.note.\textsc{pl} fake with which they.played\\
`the fake money with which they used to play’ (only: `counterfeited money’) \hfill \citep[21]{cinque:2014:semantics}

In prenominal position -- remember the position only having a direct modifier reading in Italian -- \textit{falso} has a reading as `\textit{Monopoly money}’, whereas in postnominal position it only has a reading as `\textit{forged/counterfeited money}’. As expected, this reading corresponds to the reading in predicative position.

\exg. Quelle banconote sono false.\\
that.\textsc{pl} bank.note.\textsc{pl} are fake\\
1. $\approx$ `This money is counterfeited.’\\
2. $\neq$ `This money is Monopoly money.’ \null\hfill \citep[21]{cinque:2014:semantics}

\subsubsection{Adnominal degree modifiers}

A final group is what \citet{morzycki:2016:mod} calls \is{adnominal degree modifier}\textit{adnominal degree modifiers}.

\ex. He is a complete / total / utter / absolute / true idiot. \\ \null\hfill \citep[56–57]{morzycki:2016:mod}

As the \is{adnominal degree modifier}name suggests, \citet{morzycki:2014:where, morzycki:2016:mod} does not see these elements as lexical adjectives, but rather as functional modifiers that denote a (often, but not always extreme) point on a scale for certain nouns. Nevertheless, they are relevant for this study since they cannot be used predicatively.

\ex. *This idiot is complete / total / utter / absolute / true.

\subsection{Relational adjectives} \label{RA}

\subsubsection{Nominal derivation and syntactic constraints}

\textit{Relational adjectives}, \is{relational modifier}such as \textit{environmental} or \textit{technical}, make up a relatively rich research field. They have been discussed predominantly with regard to Romance languages and English, but also in Japanese \citep{bisetto:2010:RA, nagano:2015:RA, nagano:2016:are, nagano:2018:conversion, shimada:2018:relational, morita:2010:internal, morita:2011:three, morita:2013:morph} and will be discussed in Section \Ref{RAJapanese}. Morphologically, they are characterized via their denominal appearance, exhibiting typical suffixes such as -\textit{al}, \textit{ic}(\textit{al}) in English or -\textit{ar}, -\textit{lich} or -\textit{isch} in German. This oftentimes gives relational adjectives a scientific flavor (cf. \citealt{gunkel:2008:constraints}).\footnote{They are therefore also sometimes called ``denominal adjectives" or ``(denominal) pseudo-adjectives" \citep{postal:1969:anaphoric, bartning:1980:remarques}. See \citet[52–56]{bortolotto:2016:RA} for an overview.} Semantically, they denote a specific relationship to the nouns they modify and do not just simply add a quality, as would ordinary \textit{qualitative} adjectives. Syntactically, the following cross-linguistic characteristics of relational adjectives have been attested:\footnote{For English and Romance see \citet{postal:1969:anaphoric, downing:1977:creation, levi:1978:syntax, bartning:1980:remarques, bartning:1984:aspects, bartning:1986:paralelisme, warren:1984:classifying, bosque:1993:sobre, bosque:1996:postnominal, mcnally:2004:relational, fabregas:2007:internal, fabregas:2014:adjectival, gunkel:2008:constraints, roy:2010:deadjectival, zifonun:2011:RA, rainer:2013:can, bortolotto:2016:RA}, and for an overview \citet[§ 3.3 and 9.2]{alex:2007:nounphrase}. They have also been discussed in German \citep{hotzen:1967:gegenwart, schaublin:1972:probleme, eichinger:1979:syntaktische, frevel:2005:RA, trost:2006:das}, Russian \citep{mezhevich:2002:english}, Polish \citep{cetnarowska:2015:categorial}, Hungarian \citep{kenesei:2014:on} and Selkup \citep{spencer:2017:denominal, nikolaeva:2020:mixed}. See \citet[44]{nagano:2016:are} or \citet[§ 2.3]{bortolotto:2016:RA} for similar overviews.} \is{English} \is{Romance} \is{German} \is{Russian} \is{Selkup} \is{Hungarian} \is{Polish}

\ex. Characteristics of \is{relational modifier}relational adjectives
\a. They cannot be used predicatively. (*this crisis is environmental)
\b. They are not gradable. (*more/most environmental crisis; *very environmental crisis)
\c. They have to be strictly adjacent to their head noun and cannot be combined with regular qualitative adjectives. (*environmental huge crisis; *huge and environmental crisis; *environmental and huge crisis)
\d. They appear in certain positions from which they cannot move, typically they occur in postnominal position in Romance languages. (*\textit{chimico processo}; \textit{processo chimico} `chemical process', Italian; \citealt[59]{bortolotto:2016:RA})
\e. They cannot be nominalized. (*environmental-ity)
\f. They lack polarity/do not come in antonym pairs. (*un-environmen\-tal; *non-environmental)

The fact that they cannot be used predicatively defines them as direct modifiers as tested in many different \is{relational modifier}languages. \is{English} \is{German} \is{French} \is{Russian}

\ex. \a. *My appointment this afternoon is dental. (English) \null \hfill \citep[49]{rae:2010:ordering}
\b. *This linguist is theoretical.

\exg. *Dieser Linguist ist theoretisch. (German)\\
this linguist is theoretical \\

\exg. *L’élection est présidentielle. (French)\\
\textsc{det}.election is presidential \\\null \hfill \citep[68]{bortolotto:2016:RA}

\exg. *Magazin byl kni\v{z}-n-yj. (Russian)\\
store was book-\textsc{ra-nom.sg.m} \\ \null \hfill \citep[100]{mezhevich:2002:english}

Importantly, relational adjectives can occur in predicative position in certain circumstances. One involves the class shift into a qualitative adjective (cf. \citealt{nagano:2018:conversion}).

\ex. The way he approached this problem was very theoretical.

This shift can be facilitated by degree adverbs such as \textit{very} as in \ref{italian party}. The \is{ethnic modifier}\textit{ethnic} adjective \textit{Italian} (\citealt{alex:2011:ea}, see below) has acquired a qualitative meaning as shown by the available paraphrases.

\ex. \label{italian party} \a. The party yesterday was very Italian.
\b. $\neq$ `The party took place in Italy, was organized by Italians, etc.’ (relational reading)
\c. $\approx$ `The party featured many Italian dishes, music, etc.’ (qualitative reading)

Another way in which relational adjectives can be used predicatively is when they have a \is{kind-reading}\textit{kind-reading} \citep{mcnally:2004:relational} as in the following example from \is{Catalan}Catalan.

\exg. La tuberculosi pot ser pulmonar. (Catalan)\\
\textsc{det} tuberculosis can be pulmonary\\
`The tuberculosis could be pulmonary.’ \\ \null \hfill \citep[189]{mcnally:2004:relational}

This example essentially allows the \is{relational modifier}paraphrase ``This tuberculosis can be of the pulmonary kind/type". From a syntactic point of view, a noun such as \textit{kind} or \textit{type} seems to be involved and is recovered.

Relational adjectives also resist \is{gradability}gradability, at least with \textit{indefinite degree adverbs} \citep{fabregas:2014:adjectival} of the type \textit{very}. Many are however compatible with so-called \is{modifiers of extension}\textit{modifiers of extension} (\citealt{fabregas:2014:adjectival}; called `\textit{endpoint-oriented modifiers}' in \citealt{kennedy:2005:scale}) such as \textit{completely} or \textit{100 percent}.

\ex. This is a *very / completely / 100 percent environmental problem.

Relational modifiers, typical for direct modifiers, are not only very low in the DP, they are also lower than many other direct modifiers, predicting strict adjacency to the noun. This holds cross-linguistically as shown for French below. Relational adjectives, such as \textit{industriel} `industrial’, obligatorily appear postnominally, and must be closer to the noun than qualitative adjectives, whatever status one assumes for them. \is{French}

\ex. \ag. un développement industriel rapide (French)\\
a development industrial rapid\\
`a fast industrial development’
\bg. *un développement rapide industriel\\
a development rapid industrial \null \hfill \citep[61]{bortolotto:2016:RA}\\

In Japanese, relational adjectives have been identified amongst the group of \textit{no}-modifiers, free of morphological processes (\citealt{bisetto:2010:RA, nagano:2016:are} a.o.). Importantly they do fulfill the relevant characteristics, a restriction on predicativity and gradability.

\ex. \ag. kankyō-no kiki \\
environment-\textsc{no} crisis\\
`environmental crisis’
\bg. *Kono kiki-ga kankyō-da.\\
this crisis-\textsc{nom} environment-\textsc{cop}\\
(`This crisis is environmental.’)
\cg. *totemo kankyō-no kiki \\
very environment-\textsc{no} crisis\\
(`very environmental crisis’)

This suggests \is{relational modifier}that those lexemes are simply nouns, since no obvious derivational patterns are detectable. Furthermore, relational adjectives in this language do not compete with other constructions, such as compounds or PPs with the same meaning, which is very well observed in European languages (for example \textit{nominal modification} vs. \textit{noun modification} in \is{English} English, \citealt{hacken:2013:compounds, rainer:2013:can, nikolaeva:2020:mixed}). Again, Japanese has only one option of modification at its disposal where other languages can show variation.

\subsubsection{Classificatory, thematic and ethnic adjectives}

While in many cases, relational adjectives classify the type of noun they modify, they can also express argument status. This happens when they modify event or action nouns \citep{grimshaw:1991:argument}.

\ex. environmental protection

Here, \textit{environmental} expresses the role of the internal argument of the verb \textit{to protect} underlying the action noun.\footnote{Theoretically, it could express a subject role, but our world knowledge tells us that the environment needs to be protected, rather than is the protector.} Relational adjectives used this way have been called \textit{thematic} since they ``saturate semantic roles" \citep[360]{bosque:1996:postnominal} contrasting them with \textit{classificatory relational adjectives}. Thematic relational adjectives are interesting since cross-linguistic variation in the type of argument they can express has been detected. For instance, in \is{German} German, object readings of thematic adjectives are categorically impossible \citep{gunkel:2008:constraints}.

\exg. städt-isch-e Reinigung (German)\\
city-\textsc{ra-nom.f.sg} cleaning.\textsc{nom.f.sg}\\
1. $\approx$ `cleaning by the city' (subject) \\
2. $\neq$ `cleaning of the city' (object) \\ \null \hfill \citep[294]{gunkel:2008:constraints}

Another relevant subgroup includes so-called \is{ethnic modifier}\textit{ethnic adjectives}. Ethnic adjectives are toponymic modifiers which denote a geographical origin and sometimes even titles and can be thematic or classificatory \citep{giorgi:1991:syntax, alex:2011:ea, arseni:2014:EA, kotowski:2016:adj}. Arguably, they also show restrictions in their use as thematic modifiers, resisting the role of direct objects. This goes back to the following classic example of \citet[111]{kayne:1981:ecp}.

\ex. \a. Everyone deplored the Russian destruction of China. (Subject reading)
\b. *Everyone deplored the Chinese destruction by Russia. (Object reading)

On the other hand, the \textsc{theme} reading is possible in English with genitive modifiers \citep{kayne:1981:ecp}.

\ex. Everyone deplored China's destruction by Russia. (Object reading)

Classificatory ethnic adjectives exist as well. As they correspond to the categories \textsc{nationality}/\textsc{origin}/\textsc{provenance} put forth in \citet{sproat:1991:the, scott:2002:stacked, laenzlinger:2005:french, cinque:2010:adj}, they are often not treated as true ethnic adjectives. Importantly, contrary to \is{relational modifier}relational adjectives, these adjectives can, at least in some languages, appear in predicative position, including \is{English} English \ref{ethnic english} and \is{Greek} Greek \ref{ethnic greek}.

\ex. These bed sheets are Italian. \label{ethnic english}

\exg. I tsanta tis ine italiki. (Greek)\\
\textsc{det} bag \textsc{poss.3.sg} \textsc{cop} Italian\\
`Her bag is Italian.' \null \hfill \citep[138]{alex:2011:ea} \label{ethnic greek}

In Japanese, modifiers that express \textsc{nationality} only appear among the group of \textit{no}-modifiers and can never be used predicatively, neither with object nor with event nouns. This highlights their status as direct-only modifiers.

\ex. \ag. *Kono kuruma-wa Nihon-da. \label{nihon}\\
this car-\textsc{top} Japan-\textsc{cop}\\
(`This car is Japanese.’)
\bg. *Kono kōgeki-wa Nihon-da.\\
this attack-\textsc{top} Japan-\textsc{cop}\\
(`This attack is Japanese.’)

\subsubsection{Categories and ordering} \label{ordering relational}

Relational adjectives also show \is{adjective ordering restrictions}ordering restrictions among each other analyzed in \citet{bortolotto:2016:RA} for many Romance languages. \is{Italian}

\ex. \ag. escursioni montane estive (Italian)\\
hike.\textsc{pl} mountain.\textsc{ra.pl}. summer.\textsc{ra.pl}\\
\bg. ??escursioni estive montane\\
hike.\textsc{pl} summer.\textsc{ra.pl} mountain.\textsc{ra.pl}\\
`summer mountain hikes’ \null\hfill \citep[148]{bortolotto:2016:RA}

The systematic \is{relational modifier}hierarchy established in \citet{bortolotto:2016:RA} reads as \is{adjective ordering restrictions}follows.

\ex. \emph{Hierarchy of relational adjectives (based on \citealt[180]{bortolotto:2016:RA}})\\
\textsc{time} $>$ \textsc{location} $>$ \textsc{agent} $>$ \textsc{instrument} $>$ \textsc{matter} ($>$) \textsc{theme} $>$ N

\subsection{Summary}

Both non-intersective and relational modifiers are thus groups of modifiers whose existence in Japanese immediately verifies the existence of direct modification.

\section{Indirect modification} \label{dP}

This section outlines the two types of indirect modifiers and briefly discusses the syntax of indirect modifiers focusing on the derivation adopted in this book. Indirect modification will be the topic of Chapter \ref{indirect japanese}.

\subsection{Two kinds of indirect modifiers: Full vs. reduced relative clauses} \label{chapter RC}

\subsubsection{General differences}

As mentioned above, indirect modifiers come in two types: \textit{Full relative clauses} and \is{reduced relative clause}\textit{reduced relative clauses} (RRCs) \citep{cinque:2010:adj, cinque:2020:relcl, krause:2001:rrc, douglas:2016:syntactic, harwood:2018:rrc}. Traditionally, reduced relative clauses have been assumed to be derived from full relative clauses \citep{smith:1964:det} and the following characteristics are often invoked that highlight their reduced/impoverished nature (\citealt[203]{cinque:2020:relcl}; \citealt{bhatt:1999:covert, krause:2001:rrc, douglas:2016:syntactic, harwood:2018:rrc}).

\ex. Characteristics that differentiate reduced from full relative clauses \label{chara rrc}
\a. They do not make use of relative pronouns or complementizers.
\b. Fewer positions on the Accessibility Hierarchy \citep{keenan:1977:noun} can be relativized/only the subject can be relativized.
\c. Tense is reduced/No tense is involved.
\d. No overt subject appears.
\e. A restriction to certain types of verbs applies.

These characteristics seem accurate in languages such as \is{English} English. But even there, problems arise when trying to differentiate \is{reduced relative clause}reduced relative clauses from adjectives, if they do not appear postnominally or are participial or infinitival. Testing the status of a relative clause in a prenominal-RC language,\footnote{With this notion I refer to languages with only prenominal relative clauses in line with \citet{wu:2011:syntax}.} however, is almost impossible based on these criteria alone. That is mainly because in such languages, usually only one surface form for relative clauses exists and relative pronouns/complementizers are typically absent. While these facts speak for a reduced status, the fact that overt subjects typically appear freely and many positions on the Accessibility Hierarchy can be relativized speak for a larger status (cf. \citealt{wu:2011:syntax}).

In the cartographic program, two structural differences have been identified for full and reduced relative clauses: 1. The internal make-up of reduced relative clauses contains less functional material; 2. They occupy different merging positions.

Regarding the first point, while full relative clauses are of the structural size CP, reduced relative clauses are (arguably) smaller, although their exact size is controversial. In seemingly well-known languages such as English, convincing evidence, mostly in the form of co-occurrence with high adverbials \citep{douglas:2016:syntactic, harwood:2018:rrc, cinque:2020:relcl} can be found that they must be at least of size IP. On the other hand, some authors have argued that reduced relatives in prenominal-RC languages are only vPs \citep{krause:2001:rrc}, while some assume a specific PrtP projection for participles \citep{bhatt:1999:covert}.

\subsubsection{Ordering and the syntax of the indirect domain}

Ordering restrictions, arguably, provide a better testing ground. According to \citet{cinque:2010:adj, cinque:2020:relcl}, the indirect domain (the area up to dP, see below) looks as depicted in the (simplified) tree in \figref{Indirect domain Cinque}.

\begin{figure}
\centering
 \begin{tikzpicture}[baseline=0pt, remember picture]
\tikzset{level distance=25pt, sibling distance=1pt}
 \Tree [.DP {} [.FP\textsubscript{RC} CP [.F′\textsubscript{RC} F [.FP\textsubscript{Num} NumP [.FP\textsubscript{RRC} IP [.F′\textsubscript{RRC} F [.dP d [.FP\textsubscript{direct} {} [.... {} [.NP \edge[roof]; \node (NP) {N}; ] ] ]]] ]]]] ]
 \end{tikzpicture}
\caption{Structure of the indirect domain according to \citet{cinque:2010:adj, cinque:2020:relcl}}
\label{Indirect domain Cinque}
\end{figure}

Within the indirect domain, the projections for full relative clauses are higher than the projections for reduced relative clauses. Recall that several projection for the same type have to be assumed, explaining cases of relative clause stacking. Furthermore, there is a projection for numerals intervening. This predicts the following overall hierarchy.

\ex. (\textsc{demonstrative}) ($>$) \textsc{relative clause} $>$ \textsc{numeral} $>$ \textsc{reduced relative clause} ($>$) (dP) ($>$) (\textsc{direct adjective}) ($>$) (N)

In fact, the interaction with \is{numeral}numerals serves as an important test bed as well: Languages with a homomorphic order of constituents, that is where modifiers appear on the same side of the noun, show the cross-linguistically stable tendency that full relative clauses precede numerals in linear order, while \is{reduced relative clause}reduced relative clauses follow them. The tendency that full relative clauses typically precede numerals has been identified already in \citet{greenberg:1963:some} and \citet{hawkins:1983:word} and is further discussed in \citet{cinque:2005:deriving, cinque:2023:linearization} and especially \citet[§ 3.5]{cinque:2020:relcl}.\footnote{Similarly, relative clauses precede adjectives, if they are on the same side of the noun, see \citet{cinque:2005:deriving, cinque:2023:linearization} and \citet[298–299]{rijkhoff:2002:noun}.} This is shown in the examples below, which are taking from a rigid head-final and head-initial language each. Here, the hierarchy \textit{demonstrative} $>$ \textit{relative clause} $>$ \textit{numeral} is observable (cf. \citealt[225]{cinque:2020:relcl}). \is{Wolaytta} \is{Lango}

\exg. [he taa-w kuttuwa ehida] iccashu adussa laagge-t-i\\
those I-\textsc{dat} chicken having.brought five tall friend-\textsc{pl-nom}\\
`those five tall friends [who brought me a chicken]’ \\ \null \hfill Wolaytta, West Cushitic, head-final (\citealt[215]{lamberti:1997:wolyatta})

\exg. gwóggî à dòŋò àryó-nì [ámê lóc\textschwa {} ònèkò]-nì \\
dogs \textsc{att.part} big two-this \textsc{rel} man kill.\textsc{3.sg}.\textsc{perf}-this\\
`these two big dogs [that the man killed]’ \\ \null \hfill Lango, Eastern Nilotic, head-initial (\citealt[156]{noonan:1992:lango})

\largerpage
On the other hand, the underlying order \is{numeral}\textit{numeral} $\succ$ \textit{RRC} can be verified in \is{German} German (\ref{numeral rrc german}; cf. also \citealt{eichinger:1991:ganz, trost:2006:das}) and English \ref{numeral rrc english}. Numerals and participial \is{reduced relative clause}(reduced) relative clauses must both occur prenominally and are best in that order.

\ex. \label{numeral rrc german} \ag. die drei [dort parkenden]\textsubscript{RRC} Autos (German)\\
the three there parking.\textsc{prs.ptcp} cars\\
\bg. ??die [dort parkenden]\textsubscript{RRC} drei Autos\\
the there parking.\textsc{prs.ptcp} three cars\\
`the three cars parking over there’

\ex. \label{numeral rrc english} \a. these (other) two [recently completed]\textsubscript{RRC} plays
\b. ?these (other) [recently completed]\textsubscript{RRC} two plays \null \hfill \citep[229]{cinque:2020:relcl}

\subsubsection{A Note on the indefinite head d} \label{indefinitedpdiscussion}

According to \citet{cinque:2010:adj, cinque:2020:relcl}, the threshold between the direct and the indirect domain is the indefinite head \is{dP}d, which is the ``head of the relative clause itself" \citep[34]{cinque:2010:adj} and bears an identical constituent within the relative clause. The syntactic role of d as the \textit{external head} and counterpart to the \textit{internal head} is essentially relevant in Cinque’s double-headed relative clause theory \citep{cinque:2020:relcl}. From a semantic point of view, d assigns ``referential import" \citep[34]{cinque:2010:adj} meaning that modifiers above this head modify nouns with referentiality of their own. The author seems to have in mind a more referential status of relative clauses, which, for instance, matches the intersective interpretation, whereas direct modifiers not necessarily bear this feature of referentiality, which explains their non-intersective interpretation. In this book, dP will be displayed as the threshold between the two domains, but it does not serve a syntactic function beyond that, as the double-headed structure will not be adopted for Japanese. \is{double-headed relative clause}

\subsection{The syntactic derivation of relative clauses}

\subsubsection{The underlying structure}

A plethora of different analyses can be found for the structure and derivation of (the different types of) relative clauses in English alone \citep{schachter:1973:focus, vergnaud:1974:french, vergnaud:1985:dep, kayne:1994:anti, bianchi:1999:consequences, bhatt:2002:raising, salzmann:2019:new}. What is common to them is that they must find a way to solve the \textit{Connectivity Problem} \citep{salzmann:2017:reconstruction}, i.e. the question how the correct interpretation of the noun, both inside the relative clause and inside the matrix clause, can be ascertained. Prior to derivation, I assume the underlying structure of a DP containing a full relative clause as in \figref{Structure of full relative clause}, following cartographic research until \citet{cinque:2010:adj}. \is{connectivity problem}

\begin{figure}
\centering
\begin{tikzpicture}[baseline=0pt, remember picture]
\tikzset{level distance=25pt, sibling distance=1pt}
 \Tree [.DP the [.AgrP \node (Agr) {}; [.Agr′ Agr [.FP\textsubscript{RC} [.CP {} [.C′ [.C that ] [.IP [.DP \edge[roof]; {John} ] [.I′ I [.VP {} [.V′ [.V write ] {} ] ] ] ]] ] [.F′\textsubscript{RC} F [.FP\textsubscript{Num} NumP [.FP\textsubscript{RRC} IP [.F′\textsubscript{RRC} F [.dP d [.FP\textsubscript{direct} {} [.... {} [.NP \edge[roof]; \node (NP) {book}; ] ] ]]] ]]]] ]]]
 \end{tikzpicture}
\caption{Structure of DP with relative clause \textit{the book that John wrote}}
\label{Structure of full relative clause}
\end{figure}

Full relative clauses are of size CP and merge in the specifier position of available functional projections dedicated to this type of relative clause and higher than the projection for numerals and those for reduced relative clauses. The structure for a reduced relative clause looks similar except for the fact that a CP-layer is lacking.

\subsubsection{The derivation}

Three major types of analyses have been invoked for head-external relative clauses, each of them ensuring in a different way that the noun can be interpreted correctly both inside the relative clause, where it is not present on the surface, and within the matrix clause:\footnote{Even among a single stream of analysis, several differences exist as to the fine details and technicalities of the derivation. I will simplify these details and differences in the following discussion. See \citet[§ 2]{salzmann:2017:reconstruction} for an informative overview.}

\begin{enumerate}
\item The Raising analysis
\item The Matching analysis
\item Operator movement
\end{enumerate}

In the \is{raising}\textit{raising} analysis, also known as \textit{promotion} or \textit{head-internal} analysis, the noun is structurally present inside the relative clause at the beginning of the derivation. It then \textit{raises} out of the relative clause and leaves a trace behind. This guarantees the correct interpretation \citep{schachter:1973:focus, vergnaud:1974:french, vergnaud:1985:dep, kayne:1994:anti, bianchi:1999:consequences, bhatt:2002:raising}.

According to the \is{matching}\textit{matching} analysis, no movement happens. Instead, there are two instances of the noun, one inside and one outside the relative clause. The noun inside the relative clause (in some versions the complement of an operator) is deleted at PF under phonological identity with the noun in the matrix clause \citep{lees:1961:constituent, chomsky:1965:aspects, schachter:1973:focus, sauerland:1998:the, sauerland:2003:a, salzmann:2017:reconstruction, salzmann:2019:new}.

Both analyses have been proven to be valid in different languages. Essentially, the difference boils down to the fact that in raising analyses, movement of the noun takes place. This predicts the availability of different syntactic \is{reconstruction effects}reconstruction effects, such as idiomatic readings and island effects (see \citealt{cinque:2008:more, cinque:2020:relcl, salzmann:2017:reconstruction, salzmann:2019:new} for overviews and discussions).

In this book, I will adopt the standard, older and not uncontroversial, \is{operator}\textit{operator-movement} derivation for Japanese in Section \ref{operator}, according to which a relative clause operator moves from within the relative clause, for instance from the complement position of the verb in object relative clauses, to Spec, CP. This operator is co-indexed with the head noun outside the relative clause and ensures it is interpreted correctly inside the relative clause as well (cf. \citealt{chomsky:1977:on}, especially \citealt{miyamoto:2017:RC} for Japanese). This derivation is also given as a possible option for Japanese, and other prenominal relative clause languages in \citet[81]{cinque:2020:relcl}, abstracting away from the double-headed structure. For English, the derivation looks as in \figref{Operator derivation English}, including movement of the noun to a higher functional projection for word order purposes.

\begin{figure}
\fittable{
\begin{tikzpicture}[baseline=0pt, remember picture]
\tikzset{level distance=25pt, sibling distance=1pt}
 \Tree [.DP the [.AgrP \node (Agr) {}; [.Agr′ Agr [.FP\textsubscript{RC} [.CP \node (op1) {Op\textsubscript{1}}; [.C′ [.C that ] [.IP [.DP \edge[roof]; {John} ] [.I′ I [.VP {} [.V′ [.V write ] \node (op2) {$<$Op$>$}; ] ] ] ] ]] [.F′\textsubscript{RC} F [.FP\textsubscript{Num} NumP [.FP\textsubscript{RRC} IP [.F′\textsubscript{RRC} F [.dP d [.FP\textsubscript{direct} {} [.... {} [.NP \edge[roof]; \node (NP) {book\textsubscript{1}}; ] ] ]]] ]]]] ]]]
 \draw [color=blue][->] (NP) |- +(0,-2) -| (Agr);
 \draw [color=blue][->] (op2) |- +(0,-2) -| (op1);
 \end{tikzpicture}
 }
\caption{Operator derivation of English relative clauses}
\label{Operator derivation English}
\end{figure}

The difference between the \is{raising}raising analysis and the \is{operator}operator-movement analysis is thus the presence of an operator, which ensures correct interpretability; a commonality is the prediction of \is{reconstruction effects}reconstruction effects, contrary to the matching derivation.\footnote{See Section \ref{internalstructure} for why we need to assume this structure for seemingly simpler modifiers as well.}

\subsubsection{A note on the double-headed structure}

In recent times, \citet{cinque:2020:relcl} has proposed a universal \is{double-headed relative clause}double-headed structure for relative clauses that looks as displayed in \figref{Double-headed structure English} for a full relative clause.

\begin{figure}
\centering
\begin{tikzpicture}[baseline=-0pt, remember picture]
\tikzset{level distance=25pt, sibling distance=1pt}
 \Tree [.DP [.D {the} ] [.FP\textsubscript{RC} [.CP \node (N1) {}; [.C′ [.C {that} ] [.IP [.DP {John} ] [.I′ I [.VP [.V {wrote} ] [.dp2(int) d [.NP \edge[roof]; [. \node (N2) {N}; {book} ] ] ] ] ] ] ] ] [.F′\textsubscript{RC} F [.dP1(ex) d [.NP \edge[roof]; [.N {book} ]]]]]]
 \end{tikzpicture}
\caption{Double-headed structure of English full relative clauses}
\label{Double-headed structure English}
\end{figure}

Crucial to this idea is that the noun is structurally present twice, once inside and once outside the relative clause,\footnote{In fact, what is structurally present twice, is not only the noun but any projection lower than the indefinite head d.} and deletable in either a \is{raising}raising or a \is{matching}matching fashion. Either the internal head \is{dP}dP2 moves to Spec, CP, where it receives operator status and licenses subsequent deletion of the external head (\textit{raising}), or the external head dP1 stays in situ and licenses the deletion of the internal head dP2 (\textit{matching}), cf. \citet[25–45]{cinque:2020:relcl}. This account is argued to be universal and both derivations are accessible in all languages at all times.

The problem with this account is that evidence for a \is{double-headed relative clause}double-headedness of the noun has only been proven in a small number of languages, including \is{Kombai} Kombai (Awyu, West Papua; cf. \citealt{vries:2020:greater}) and \is{Wenzhounese} Wenzhounese (Sinitic; cf. \citealt{hu:2018:new}). In most languages, including \is{English} English and Japanese (although see \citealt{erlewine:2016:unifying} for possible exceptions), there is only one overt noun, or some version of it. Furthermore, as shown in Section \ref{operator}, Japanese shows ample evidence for movement in the form of \is{reconstruction effects}reconstruction effects, which in turn speaks in favor of either an operator movement or \is{raising}raising derivation. For these reasons, I will not adopt the double-headed structure here, but a simpler structure, and assume operator movement as the operation that derives relative clauses in Japanese.

\subsection{Summary: Indirect modification}

To summarize, indirect modifiers come in two types, full and reduced relative clauses. These modifiers differ in their structural make-up and syntactic height. In this book, relative clauses will be assumed to be derived via movement of an operator from within the relative clause to a Spec, CP, which is co-indexed with the noun and ensures correct interpretability.

\section{Additional accounts}

In this section, I will present ideas from three additional proposals, that differ from the mainline approach in \citet{cinque:2005:deriving, cinque:2017:micro, cinque:2010:adj, cinque:2020:relcl} in different ways and will all enrich the system defended here. They all share the assumption of a large number and universal character of functional projections within the DP, but add different ingredients to be adopted in the following. I will focus on what is relevant in these approaches for the purpose of this study.

\subsection{Svenonius: Gradability and intersectivity} \label{svenonius}

\citet[249–250]{larson:2021:rethinking}, in a general critique on Cartography (and see \citealt{truswell:2004:attributive, truswell:2009:attributive, svenonius:2008:position, kotowski:2016:adj} and \citealt{kim:2019:syntax} for different languages and categories), mentions the \is{problem of rigidity}`Problem of Rigidity' according to which such a strict structural selection that we see in Cartography should always be reflected in linear order (\citealt{kayne:1994:anti, cinque:2010:adj} etc.) and thus deviations should never occur. However, as \citet{svenonius:2008:position} points out, sometimes the \is{adjective ordering restrictions}ordering restrictions proposed in \citet{scott:2002:stacked} and similar work seem rather ``tailored" \citep[37]{svenonius:2008:position}. For instance, the order of the modifiers of \textsc{dimension} and \textsc{shape} in \ref{longthick} is rather free. On the other hand, he does point out clear cases where ordering restrictions hold, for example between modifiers of \textsc{dimension} and \textsc{color}, illustrated in \ref{tinyred}. \is{intersectivity}

\ex. \label{longthick} \a. a long thick rope (\textsc{dimension} $\succ$ \textsc{shape})
\b. a thick long rope (\textsc{shape} $\succ$ \textsc{dimension}) \null \hfill\citep[37]{svenonius:2008:position}

\ex. \label{tinyred} \a. tiny red hats (\textsc{dimension} $\succ$ \textsc{color})
\b. ??red tiny hats (\textsc{color} $\succ$ \textsc{dimension} \null \hfill\citep[36]{svenonius:2008:position}

Svenonius dismisses \is{adjective ordering restrictions}ordering restrictions based on fine-grained semantic categories, and instead argues that such ordering restrictions only arise with modifiers with different levels of \is{gradability}gradability (as well as intersectivity). For example, modifiers of \textsc{color}, \textsc{origin} and \textsc{material} are clearly non-gradable, while modifiers of \textsc{dimension}, \textsc{quality} and \textsc{shape} are gradable (countable).

\ex. \a. ??very red / Japanese / metal hat (\textsc{color}, \textsc{origin}, \textsc{material})
\b. very big / nice / round hat (\textsc{dimension}, \textsc{quality}, \textsc{shape})

The different types of modifiers can, according to Svenonius, also be differentiated according to their level of intersectivity, and this can be tested via paraphrasing them with a \textit{for-PP}. Only subsective modifiers, for instance of \textsc{dimension}, allow such a paraphrase, intersective modifiers, for example of \textsc{color}, do not (cf. \citealt{mckinney:2011:adj} and \citealt{cinque:2014:semantics}). \is{intersectivity}

\ex. \a. big for a butterfly \null \hfill subsective
\b. *red for a hat \null \hfill intersective

\citet[§ 4.3]{svenonius:2008:position} proposes that the DP incorporates the functional projections nP, reserved for non-gradable intersective modifiers, and SortP for gradable, subsective modifiers. Since non-gradable intersective modifiers, for instance of \textsc{color}, seem to be closer to the noun than gradable subsective modifiers of \textsc{dimension}, SortP hierarchically precedes nP. This yields the following hierarchy.\footnote{PlP is a projection hosting plural morphemes and UnitP (also ClP; \citealt{svenonius:2022:span}) hosts numerals. This hierarchy mirrors the types and ordering principles of different classifiers found in classifier-languages, an idea taken from \citet{borer:2005:in}, see also \citet{rijkhoff:2002:noun} and \citet{zamparelli:2000:layers}.}

\ex. UnitP $>$ PlP $>$ SortP $>$ nP $>$ NP \null \hfill \citep{svenonius:2008:position}

Svenonius does not make any reference to the difference between direct and indirect modifiers, so what exactly merges in SortP and UnitP is neither limited to lexemes of certain word classes nor of specific modification types. Instead, modifiers merge in ``whatever position makes sense for their interpretation" \citep[39]{svenonius:2008:position}. This is undesirable for our aim of deriving a DP structure that incorporates the aspect of modification type. Furthermore, a concrete problem is that \is{adjective ordering restrictions}ordering restrictions among categories of the same degree of intersectivity, for example \textsc{color} and \textsc{nationality}, are not predicted, although we do see these even in English as also pointed out in \citet{panayidou:2013:in}.

\ex. ??Chinese blue vase

The same holds true for a recent proposal extended in \citet{manzini:2025:left}, who argues in a similar vein that ordering is not determined directly in the syntax, and never based on fine-grained functional projections, but is the result of externalization \citep{chomsky:2013:problems}. Nevertheless, as she proposes, adjectives merge at different levels of the DP, depending on whether they are \is{intersectivity}gradable. Concretely, \citet[10–16]{manzini:2025:left} proposes that \is{gradability}gradable adjectives merge at the edge of the nP phase, assuming that the DP can be divided into several phases, while non-gradable adjectives merge as complements of nP.\footnote{\citet{manzini:2025:left} does not make the same reference to intersectivity concretely since, as she argues, intersectivity and subsectivity are semantic concepts and, according to her, are difficult to argue to have syntactic dimensions. Furthermore, she provides the data point in \Next where an intersective adjective precedes a subsective adjective in English, contrary to expectation. It should be noted, however, that different mechanisms might be at play here, such as scope dependencies, similar to cases addressed in
Section \ref{optionality} for Japanese.

\par

\ex. It is not true that there are no beautiful good linguists. \null \hfill \citep[10]{manzini:2025:left}

\par
}
However, data points such as those provided above are equally problematic for her account.

Besides the idea that fine-grained functional projections might not (always) be at stake, what will also be adopted from Svenonius is a specific projection on the right edge of the DP for \is{idiomatic modifier}idiomatic modifiers. He assumes that such modifiers are embedded in category-less root phrases \citep{marantz:1997:no}, importantly, below the nP level. That they occupy the position right above the noun becomes evident, since nothing can intervene between idiomatic modifiers and the noun they modify without changing the idiomatic meaning \citep[36–38]{svenonius:2008:position}.

\ex. \a. Minnesotan wild rice
\b. wild Minnesotan rice (only `uncultivated rice from Minnesota’) \\ \null \hfill\citep[36–37]{svenonius:2008:position}

\subsection{Laenzlinger: Split-DP} \label{laenzdp}

Christopher Laenzlinger (\citeyear{laenzlinger:1998:comparative, laenzlinger:2005:french, laenzlinger:2010:cp, laenzlinger:2011:elements, laenzlinger:2015:the}) focuses on the parametric differences between Romance (French) and Germanic and extends his findings to a cross-linguistic perspective. He presents a study which discusses the CP and DP in 14 typologically diverse languages, one of these languages being Japanese. His approach and DP-structure are essentially similar to Cinque’s modulo certain differences.

What is crucial for the account presented in this study is his adoption of a DP-external domain, a `\textit{Vorfeld}’ or `ω-domain’ as he calls it, above the `middle field’ (and `low field’) (\citealt[7]{laenzlinger:2011:elements}; see \citealt{grohmann:2003:prolific} for a similar account). This area can be reached via individual movement of modifiers and is headed by a dedicated D-projection DP\textsubscript{determination} above the \textit{internal} DP\textsubscript{deixis} \citep[668–669]{laenzlinger:2005:french}. The availability of two D-positions accounts for the overt occurrence of multiple D-elements in certain languages (including \is{Lakota} Lakota and \is{Urama} Urama, see \citealt[187–193]{becker:2018:articles}) as well as well-known phenomena such as \is{Romance} Romance superlatives \citep{bruge:2002:positions, bernstein:2001:focusing} or \is{Greek} Greek determiner spreading (\citealt{androutsopoulou:1995:ds}, but see \citealt{lekakou:2012:poly} for an alternative claim). Furthermore, the necessity of a left periphery in the nominal realm has been put forth by various authors to account for high topic and focus positions and movement to those (\citealt{giusti:1996:is, dimitrova:1998:fragments, ihsane:2001:specific, aboh:2004:topic, aboh:2006:when, grohmann:2003:prolific, grohmann:2015:demonstrative}; cf. \citealt[173–178]{laenzlinger:2011:elements}). In DPs with only one determiner, this determiner always starts in the lower DP but can move up in case material has moved above it \citep[667]{laenzlinger:2005:french}.

The necessity of a \is{left periphery}left periphery as a landing site for modifiers serves mainly the purpose of accounting for otherwise inexplicable placement of modifiers. As Laenzlinger assumes that the canonical position of adjectives in French is postnominal, which is the result of NP-movement to satisfy agreement between adjective and noun, \is{pied-piping}pied-piping the adjective in the process \citep[657–662]{laenzlinger:2005:french}, prenominal adjectives are only possible in the following circumstances which are all related to deictic or pragmatic contributions \citep[666]{laenzlinger:2005:french}. This echoes the arguments from the other authors mentioned above who proposed a left periphery to account for high modifiers with a topic or \is{focus}focus interpretation. \is{French}

\ex. \a.They have an emphatic or subjective interpretation such as \textit{superbe} `superb’.
\b. They are quantifiers.
\c. They appear in weak forms (\textit{gros} in the literal sense of `big’ as opposed to `grand’).

Such modifiers move to the specifier position of dedicated projections FocP/SubjectiveP, for \is{focus}focused/subjective adjectives, QuantP, for quantifiers, and WeakP for weak forms of \is{French} French adjectives respectively \citep[666–668]{laenzlinger:2005:french}. In a DP such as \ref{laenzlinger example}, first the NP has moved cyclically to an agreement position in the DP-external domain, where it checks agreement with the adjectives \textit{rouges} `red’, \textit{magnifiques} `beautiful‘ and \textit{petites} `small’. The adjective \textit{petites} subsequently moves to WeakP in the left periphery for pragmatic reasons \citep[667]{laenzlinger:2005:french}.

\exg. de petites voitures rouges magnifiques \\
\textsc{det} small car red beautiful\\
`beautiful small red cars’ \null \hfill \citep[668]{laenzlinger:2005:french} \label{laenzlinger example}

The DP-structure, given in hierarchical format, proposed by \citet{laenzlinger:2005:french, laenzlinger:2011:elements} reads as follows.

\ex. DP\textsubscript{external/deictic} $>$ QuantP $>$ FocP $>$ WeakP $>$ FP\textsubscript{AgrNP} $>$ DP\textsubscript{internal/determination} $>$ FP\textsubscript{PP} $>$ FP\textsubscript{AgrNP} $>$ FP\textsubscript{Pred} $>$ FP\textsubscript{Adj} $>$ nP $>$ NP \null \hfill \citep[35]{laenzlinger:2011:elements}

Notably, Laenzlinger does not present a systematic difference between direct and indirect modification and his projection PredP (FP\textsubscript{Pred}) is reserved for what he calls ``external adjectival predicates” \citep[670]{laenzlinger:2005:french} explaining cases of three postnominal adjectives in French. This scenario is atypical, but possible, if the final one receives an emphatic reading such as in \textit{une voiture rouge italienne MAGNIFIQUE} `a BEAUTIFUL red Italian car’ \citep[663–664, 671–672]{laenzlinger:2005:french}.\footnote{Also, the question remains why languages should exist that have agreement, but where prenominal adjectives are inflected, and postnominal adjectives uninflected, a point also raised in \citet[20, Footnote 9]{bortolotto:2016:RA}. A good example is German, where inflected attributive adjectives in prenominal position are the standard case, but postnominal adjectives are possible, in which case they are not inflected (cf. \citealt{trost:2011:whiskey}) as in \textit{Whisky pur} `pure whiskey’ \citep[383]{trost:2006:die}.} \is{French} \is{German}

The main point to be adapted from Laenzlinger is the \is{left periphery}left periphery within the DP and the possibility for individual movement. Specifically, Japanese also freely exhibits different kinds of modifiers in a position above demonstratives, among them quantifiers and direct modifiers, but no pragmatic factors are detectable.

\ex. \ag. takusan-no ano hon\\
a.lot-\textsc{no} that book\\
`those many books’
\bg. izen-no kono rūru\\
former-\textsc{no} this rule\\
`this former rule’

\subsection{Kim: A tripartite DP-structure based on Korean} \label{kimsaccount}

\citet{kim:2019:syntax} blends and amends the two preceding analyses of \citet{svenonius:2008:position} and \citet{laenzlinger:2005:french} and enriches them with her own ideas drawing from \is{Korean} Korean data. Essential to her account as well is a tripartite DP-structure consisting of three layers: A \textit{High Field} headed by DP\textsubscript{d(eixis)/r(eferential)}, a \textit{Middle Field} headed by DP\textsubscript{q(uantification)} and a \textit{Low Field} headed by DP\textsubscript{p(redication)}. These DPs license [deictic, referential, definite], [quantificational] and [predicative] features respectively \citep[123]{kim:2019:syntax}.\footnote{These DPs correspond to the layered structural mapping of satellites around a nominal or verbal nucleus proposed in \citet[218]{rijkhoff:2002:noun}, and also to three different types of nominals, namely referential, quantificational and predicative \citep[114]{kim:2019:syntax}.} Furthermore, she assumes independent movement of modifiers for reasons of feature-checking (and potentially at LF for semantic/pragmatic reasons) and even of whole DP domains and all material contained in them. Every functional projection, and this includes every DP projection, can multiply and ``create" additional functional projections of the same kind in an adjunct fashion.\footnote{``FPs of the same kind can be created in the \is{left periphery}left periphery of a DP structure as needed, forming an adjunction structure" \citep[121]{kim:2019:syntax}, cf. \citet{koopman:2000:verbal}.}

What will be adopted here, besides the split-DP also argued for in \citet{laenzlinger:2005:french, laenzlinger:2011:elements}, is her projection SpplP (SupplementaryP) situated above the highest DP and dedicated to \is{non-restrictive relative clause}\textit{supplementary} relative clauses, which are not properly \textit{integrated} in the DP \citep[116–119]{kim:2019:syntax}.\footnote{Note that \textit{supplementary} is not equal to \textit{non-restrictive}. While supplementary relative clauses are always non-restrictive, both non-restrictive and restrictive relative clauses can be integrated parts of the internal DP \citep[120]{kim:2019:syntax}.} She gives the following example (from \citealt{huddleston:2005:student}).

\ex. The doctor, [with no name tag on], got into serious trouble. \\ \null \hfill \citep[117]{kim:2019:syntax}

A (reduced) relative clause such as \textit{with no name tag on} is not essential to the DP, but \textit{supplementary}. It is base-generated in the left periphery outside the scope of D. In Japanese, non-restrictive, or supplementary, relative clauses also obligatorily occur above demonstratives (example adapted from \citealt[7–8]{ishizuka:2008:RC}).

\exg. Ito-san-ni-wa musuko-ga hitori i-ru. [[Sakunen isha-ni nat-ta] sono musuko-ga] kekkon-shi-ta.\\
Ito-Mr./Mrs.-\textsc{loc-top} son-\textsc{nom} one be-\textsc{prs} last.year doctor-\textsc{loc} become-\textsc{pst} that son-\textsc{nom} marriage-do-\textsc{pst}\\
`Mr./Mrs. Ito has one son. This son, who happened to become a doctor last year, got married.' \\ \null \hfill non-restrictive relative clause $\succ$ demonstrative $\succ$ N

In the \is{left periphery}lower DP domain, Kim assumes the same projections as \citet{svenonius:2008:position} and adds a projection for possessive modifiers and demonstratives called LocP (the satellite `\textit{Location}’ in \citet{rijkhoff:2002:noun}). She offers the following hierarchy.\is{non-restrictive relative clause}

\ex. {RC\textsubscript{Sppl}} $>$ DP\textsubscript{deictic/referential} $>$ DP\textsubscript{quantification} $>$ DP\textsubscript{predication} $>$ FocP $>$ LocP $>$ UnitP $>$ PlP $>$ SortP $>$ nP $>$ $\sqrt{}$P $>$ NP \null \hfill \citep[132]{kim:2019:syntax}

Another important assumption to be adopted (see Section \ref{size,indirect}) is the unification of full and reduced relative clauses to one type of relative clause of the same size.

\section{Further issues}

In this section, I raise two additional issues. One anticipates potential problems that come with Cartography, in the other, I briefly mention previous approaches to Japanese as part of this research program.

\subsection{Potential problems of Cartography} \label{problemscarto}

A system as rigid as Cartography cannot be free of problems and criticism. Anticipating the discussion, I will discuss several potential problems below.\footnote{For this section, consider the recent discussions in \citet{manzini:2025:left} and \citet{cinque:2025:superiority}.}

One objection concerns the existence of movement operations for the sake of order and devoid of syntactic and semantic constraints, often dubbed ``meaningless movement". Several researchers have argued that this kind of movement should be kept apart from narrow syntax (most notably see \citealt{chomsky:2019:generative} and references cited therein). Recall from Section \ref{japanese no movement} that no NP-movement needs to be assumed in Japanese since the noun is always the final element in the DP.\footnote{For meaningless movements also see \citet{hall:2015:spelling}.} \is{meaningless movement}

While this is a more ideological point, a more concrete point that begs further discussion in the future is the nature of the cross-linguistic restrictions on movement operations. How are they spread across different languages and language groups? Why do we only witness a particular type of movement, for instance progressive \is{pied-piping}pied-piping, in a given language, while other languages allow both types (see Footnote \ref{footnotemovement} in Section \ref{left-right})? Furthermore, experimental approaches have cast doubt on universal syntactic principles and related \is{absence of ordering restrictions}ordering principles rather to cognitive factors, such as subjectivity and information gain. See \citet{scontras:2017:sub, dyer:2024:evaluating}, the works cited therein and similar works by the same authors.

Sometimes, even specific subclasses of modifiers in a given language are restricted to specific positions and suggest a restriction to one particular type of movement. An interesting case are uninflectable adjectives of the type \textit{blu} `blue’ in \is{Italian} Italian, and perhaps in other Romance languages. Italian color adjectives, although restricted, can appear in prenominal and postnominal position, but this is never possible for uninflectable color adjectives such as \textit{blu} illustrated below, and neither for \is{relational modifier}relational adjectives in this and other Romance languages. \is{uninflectability}

\newpage

\ex. \ag. la (rossa) distesa (rossa) del mare (Italian)\\
the red.\textsc{f.sg} expanse.\textsc{f.sg} red.\textsc{f.sg} of.the sea\\
`the red expanse of the sea'
\bg. la (*blu) distesa (blu) del mare\\
the blue.\textsc{$\emptyset$} expanse.\textsc{f.sg} blue.\textsc{$\emptyset$} of.the sea\\
`the blue expanse of the sea' \null \hfill \citep[110]{mattiuzzi:2023:influential}

Such adjectives pose a problem for standard cartographic accounts, since no obvious relation between morphology and syntax can be drawn. Specifically, there is no obvious reason why these specific modifiers must be subject to \is{pied-piping}pied-piping. The problem also extends to the system of \citet{laenzlinger:2005:french, laenzlinger:2011:elements}, which would predict only agreeing adjectives to surface postnominally. For recent discussion see also \citet{mattiuzzi:2023:influential}.

Equally problematic are data points, provided by several authors from different languages, that show that direct modifiers (adjectives) adhere to the strict \is{absence of ordering restrictions}ordering principles dictated by Cartography cross-linguistically to different degrees (see for an overview \citealt{manzini:2025:left}). I refer the reader to Section \ref{discussiondirect}, where I discuss this point, specifically with regard to Japanese, but also to Persian and Modern Greek. Here, I would like to exemplify this point with recent data from \is{Naasioi} Naasioi, a Papuan language. See Section \ref{more languages} for discussion on this and another problematic language.

\citet{brown:2024:adj} provides very clear evidence that adjectives in Naasioi which are unequivocally direct do not have to adhere to ordering principles, as shown below (see the relevant paper for evidence that the provided adjectives are direct).

\ex. \ag. pangkaing mutaanu mosi (Naasioi)\\
big black dog\\
\bg. mutaanu pangkaing mosi\\
black big dog\\
`big black dog’ \null \hfill \citep[428]{brown:2024:adj}

As I will argue in this book, such cases are unproblematic if we assume that the nature of the functional projections available in the direct domain of the DP are subject to cross-linguistic variation. \is{Naasioi}

Finally, a potential problem concerns the presumably infinite number of functional heads present in the nominal, and also the clausal, domain. Recall that, if we take the cartographic idea seriously, every new category found in a certain language must be present in the functional spine of all languages. This point is also raised in \citet{larson:2021:rethinking}, who also reports the first of the two quotes below.

\ex. \a. The cartographic approach follows this idea in assuming that all languages share the same principles of phrases and clause composition and the same functional make-up of the clause and its phrases. \citep[68]{cinriz:2010:carto}
\b. Comparison of many different languages may provide evidence for determining the precise relative order of the different projections by combining the partial orders overtly manifested by different languages into what, in principle, should be a unique consistent order/hierarchy, imposed by UG. \citep[72]{cinriz:2010:carto}

A concrete problem arises if a modifier occurs that cannot readily be assigned to any of the existing categories of functional projections. An obvious solution would be to establish a new functional projection. But this would, at some point, potentially lead to an overabundance of categories, which is called the `\textit{Problem of Plenitude}’ in \citet[248–249]{larson:2021:rethinking}. This problem is circumvented in this study by assuming unspecific projections, an assumption that also solves the `\textit{Problem of Rigidity}’ (\citealt[249–250]{larson:2021:rethinking}) which concerns the \is{absence of ordering restrictions}lack of expected ordering principles between two modifiers of different categories. \is{problem of plenitude} \is{problem of rigidity}

In sum, Cartography faces the problems of meaningless movement, deviations from expected ordering restrictions, plenitude and rigidity. As its counter-advocates argue, it overgenerates and undergenerates at the same time \citep{larson:2021:rethinking, manzini:2025:left}. Still, this research topic is extremely powerful and offers the right tools and principles for an in-depth investigation of the Japanese DP. Especially important is the assumption of two different sources of modifiers and their different characteristics and derivations, furthermore theoretically predicting linear precedence. Even though there might be more lenient cases, some ordering principles still do exist. All these issues will be further addressed in Section \ref{discussiondirect}.

\subsection{Japanese and Cartography: Previous analyses and relevance}

\subsubsection{Laenzlinger on Japanese DP} \label{Laenz}

\citet{laenzlinger:2010:cp, laenzlinger:2011:elements} draws many implications for his universal DP structure from Japanese as a strict head-final language. His typological study is to my knowledge the only one that looks at the Japanese nominal domain in line with Cartography. However, other than his detailed analysis of the Japanese clause, his discussion of the Japanese noun phrase falls short. He gives the following example of three modifiers adhering to the expected \is{adjective ordering restrictions}order \textsc{quality} $\succ$ \textsc{dimension} $\succ$ \textsc{color}.

\exg. \label{chiisa}subarashi-i chiisa-na aka-i kuruma \\
wonderful-\textsc{i} small-\textsc{na} red-\textsc{i} car\\
`wonderful small red car(s)’ \hfill \citep[189]{laenzlinger:2011:elements}

He claims that the expected order is also the preferred one, but that modifiers in general can be \is{absence of ordering restrictions}permuted without a change in meaning or ``special focus effect(s)” \citep[190]{laenzlinger:2011:elements}. He concludes that ``at first sight Japanese (and also \is{Tatar} Tatar) pose a problem for previous theories to nominal modification and modifier ordering". As many researchers have done before him, he solves this problem by analyzing all adjectives as stacked relative clauses predicting possible permutation. This finds support in his analysis of \textsc{i} (and \textsc{na}) as tense morphemes (\citealt[191]{laenzlinger:2011:elements}, Footnote 106).

It was already shown that direct modification does exist in Japanese, so this analysis cannot be correct. It remains furthermore unclear why there should be a preferred order of modifiers in the first place, if permutation is always allowed, and why this order should be as in English. Another major problem is that his example \ref{chiisa} contains the attributive-only lexeme (\is{rentaishi@\textit{rentaishi}}\textit{rentaishi}) \textit{chiisa-na} `small’, which cannot function as a predicate, hence cannot be predicative and be embedded in a relative clause as shown in \ref{rentaishilaenz}. This group of modifiers is the topic of Section \Ref{rentaishi}.

\ex. \label{rentaishilaenz} \ag. chiisa-na ie\\
small-na house\\
`small house'
\bg. *Kono ie-wa chiisa-da.\\
this house-\textsc{top} small-\textsc{cop}\\
(`This house is small.’)

\subsubsection{Other Analyses}

Several studies discuss the Japanese clause from a cartographic perspective: \citet{endo:2007:locality} looks at the left periphery in Japanese, integrating Japanese into Rizzi's account of locality and Relativized Minimality (\citealt{rizzi:1997:the}). Other authors have looked at the right periphery (cf. overview in \citealt{saito:2015:carto}). \citet{ueda:2008:person} develops a hierarchy of modals expressed through sentence-final particles, \citet{saito:2013:sentence} discusses various sentence types and movement to the right periphery, and \citet{saito:2012:deriving} develop a hierarchy of complementizers and sentence-final particles which is refined in \citet{saito:2015:carto}. What is interesting about all these studies is that they do find a hierarchy and derive generalizations that are empirically and theoretically valid, but for the clause, again highlighting the more controversial and recalcitrant nature of the DP.

\subsubsection{Narrog (2009)} \label{narrog}

\citet{narrog:2009:mod} is a very detailed study on the ordering of modal elements in Japanese and, to my knowledge, the only one that combines classical approaches from Japanese linguistics and Cartography. The author identifies twelve affixes, 20 periphrastic constructions, six verbal inflections and an infinite number of modal adverbs in Japanese \citep[71–76]{narrog:2009:mod} and establishes ordering and scope principles based on a corpus study. He compares his findings to those of Cinque’s adverbial hierarchy (\citeyear{cinque:1999:adverbs}) and eventually rejects the cartographic model. Besides several mismatches,\footnote{For example, Cinque only has one position for evidentiality, Narrog observes three which all occupy different positions in relation to modal and epistemic markers.} the biggest problem, Narrog claims, is the interaction of modal categories with those of tense and aspect in his model which differs greatly from Cinque’s version \citep[233–234]{narrog:2009:mod}. In the end, Cinque’s model is deemed too abstract and ``fictional” for Narrog. He claims such models can only be maintained if they are ``understood as idealizations, or rough approximations to actual language data” \citep[243]{narrog:2009:mod}. Furthermore, he calls into question the ``psychological reality of the existence of the same universal hierarchy of categories in every sentence of any language of the world”, again pointing towards the problem of rigidity \citep{larson:2021:rethinking}. Even though he ends up rejecting Cartography as an explanation for the problems he is concerned with, his study is still relevant as it produces a very substantial hierarchy on the basis of empirical data that could in turn be checked for other languages. This hierarchy is picked up in the discussion on adverbials in Section \ref{cinque adj}. \is{modality}

\subsubsection{Takamine (2010)} \label{takamine}

\citet{takamine:2010:hierarchy} looks at the order of prepositional phrases in Japanese following the theoretical framework of \citet{cinque:1999:adverbs, cinque:2006:complement} and \citet{schweikert:2005:order}. Because Japanese is a scrambling language with a variety of possible orders, it is necessary to test the order of PPs in a neutral environment. She looks at nine different categories and employs three tests for their interaction. She arrives at the following general hierarchy (see for discussion and references \citealt[58--83]{takamine:2010:hierarchy}).

\ex. \emph {Hierarchy of PPs in Japanese (\citealt[94]{takamine:2010:hierarchy})}\\
\textsc{temporal} $>$ \textsc{locative} $>$ \textsc{comitative} $>$ \textsc{reason} $>$ \textsc{source} $>$ \textsc{goal} $>$ \textsc{instrumental}/\textsc{mean} $>$ \textsc{material} $>$ \textsc{manner}

Takamine’s study is relevant for two reasons. First, she shows that word order in Japanese is not as random as it appears at first sight. Second, this hierarchy bears some resemblance to the nominal domain. This is acknowledged in \citet{bortolotto:2016:RA}, who compares Takamine's and Rae's hierarchies (\citeyear{rae:2010:ordering}) to her hierarchy of \is{relational modifier}relational adjectives in Romance presented above.

\section{Chapter summary}

This chapter introduced the theoretical background of this study focusing on the core ideas of the cartographic account and the difference between direct and indirect modifiers. It also picked up all relevant ideas to be implemented in the account of the Japanese DP to follow. The chapter ended with a quick, condensed overview of problematic aspects of Cartography and the few studies which pursued an analysis of certain aspects of Japanese syntax under a cartographic lens, highlighting the perhaps surprising lack of studies concerned with the nominal domain.

